% !TeX program = xelatex
\documentclass[UTF8]{ctexart}

\usepackage[a4paper,margin=2.6cm]{geometry}
\usepackage{amsmath,amssymb,amsthm}
\usepackage{enumitem}
\usepackage[hidelinks]{hyperref}

\newtheorem{definition}{定义}[section]
\newtheorem{lemma}{引理}[section]
\newtheorem{proposition}{命题}[section]
\newtheorem{corollary}{推论}[section]

\title{Join Order 论文结论证明札记}
\author{}
\date{}

\begin{document}
\maketitle
\pagestyle{plain}

\section*{说明}

这份文档用于集中记录 join order 相关论文中一些关键结论的证明, 方便后续阅读和复用。

\tableofcontents

\section{MPDP}

\subsection{全局 CCP-Pair 的跨边落在同一个 block}

本小节证明 MPDP 在一般图枚举部分使用的一个关键事实：给定一个全局
CCP-Pair \((S_L,S_R)\)，所有从左侧连到右侧的边都属于同一个 block。
因此，MPDP 可以把一个全局切分的``切口''定位到某个 block 内部，再由
block-local 的 \((lb, rb)\) 通过 \textsc{grow} 扩展成全局的
\((S_L,S_R)\)。

以下约定只在本小节中使用。设 join graph 为无向图 \(G=(V,E)\)。对任意点集
\(X\subseteq V\)，\(G[X]\) 表示由 \(X\) 诱导出的子图；若 \(G[X]\) 连通，
则称 \(X\) 是一个 connected subset。在不考虑笛卡尔积的 join enumeration
中，一个 CCP-Pair \((S_L,S_R)\) 满足：
\begin{enumerate}[label=(\arabic*)]
  \item \(S_L\neq\varnothing\)，\(S_R\neq\varnothing\)；
  \item \(S_L\cap S_R=\varnothing\)；
  \item \(G[S_L]\) 和 \(G[S_R]\) 都连通；
  \item 左右两边之间至少存在一条 join graph 边。
\end{enumerate}
记左右之间的跨边集合为
\[
E_\times(S_L,S_R)=\{(u,v)\in E\mid u\in S_L,\ v\in S_R\}.
\]
一个 block 是 \(G\) 的 maximal nonseparable subgraph。这里
nonseparable 可以理解为连通且没有 cut vertex。直观地说，block 是图中无法
再被某个单点割开的最大局部区域；不同 block 之间只能通过 cut vertex 串接。

注意，\(E_\times(S_L,S_R)\neq\varnothing\) 不是本小节需要额外证明的结论。
它已经包含在 CCP-Pair 的定义中。本小节真正需要证明的是: \textbf{一旦左右之间有
一条或多条跨边，这些跨边不可能分散在多个 block 中。}

\begin{lemma}
设 \((S_L,S_R)\) 是 CCP-Pair。若
\[
e_1=(l_1,r_1),\qquad e_2=(l_2,r_2)
\]
是两条不同的跨边，其中 \(l_1,l_2\in S_L\)，\(r_1,r_2\in S_R\)，
则 \(e_1\) 与 \(e_2\) 位于同一个 cycle 上。
\end{lemma}

\begin{proof}
因为 \(G[S_L]\) 连通，所以存在一条完全位于 \(S_L\) 内部的简单路径
\[
P_L:l_1\leadsto l_2.
\]
因为 \(G[S_R]\) 连通，所以存在一条完全位于 \(S_R\) 内部的简单路径
\[
P_R:r_1\leadsto r_2.
\]

将这两条路径与两条跨边拼接：
\[
l_1 \xrightarrow{P_L} l_2
\xrightarrow{e_2} r_2
\xrightarrow{P_R^{-1}} r_1
\xrightarrow{e_1} l_1.
\]
由于 \(S_L\cap S_R=\varnothing\)，左侧路径 \(P_L\) 与右侧路径 \(P_R\)
不会共享内部顶点。取 \(P_L\) 和 \(P_R\) 为简单路径后，上述闭合 walk
包含一个同时经过 \(e_1\) 与 \(e_2\) 的 cycle。若两条跨边共享某一侧端点，
则对应一侧的路径长度可以为 \(0\)，结论仍然成立。
\end{proof}

\begin{lemma}
若两条边 \(e_1,e_2\) 位于同一个 cycle 上，则它们属于同一个 block。
\end{lemma}

\begin{proof}
设 \(C\) 是同时包含 \(e_1\) 和 \(e_2\) 的 cycle，设 \(B\) 是包含
\(e_1\) 的 block。因为 \(C\) 本身没有 cut vertex，所以 \(C\) 是
nonseparable 的；\(B\) 按定义也是 nonseparable 的。

\(B\) 与 \(C\) 至少共享边 \(e_1\)，因此至少共享 \(e_1\) 的两个端点。
对任意顶点 \(x\)，删除 \(x\) 后，\(B-x\) 仍连通，\(C-x\) 也仍连通；
并且由于 \(B\cap C\) 至少有两个顶点，删去一个顶点后二者仍至少共享一个顶点。
所以 \((B\cup C)-x\) 仍连通。也就是说，\(B\cup C\) 也是 nonseparable 的。

如果 \(C\) 中存在不属于 \(B\) 的边，那么 \(B\cup C\) 就是一个严格大于
\(B\) 的 nonseparable subgraph，这与 \(B\) 作为 maximal nonseparable
subgraph 的定义矛盾。因此 \(C\subseteq B\)，于是 \(e_2\in B\)。
\end{proof}

\begin{proposition}
设 \((S_L,S_R)\) 是一个全局 CCP-Pair。则所有跨边
\[
E_\times(S_L,S_R)=\{(u,v)\in E\mid u\in S_L,\ v\in S_R\}
\]
都属于同一个 block。
\end{proposition}

\begin{proof}
若 \(E_\times(S_L,S_R)\) 只有一条边，则结论显然成立，因为这条边属于某个
唯一的 block。

若存在至少两条跨边，任取其中两条 \(e_1,e_2\)。由上一节的路径拼接证明，
\(e_1\) 与 \(e_2\) 位于同一个 cycle 上；由 cycle 与 block 的引理，
\(e_1\) 与 \(e_2\) 属于同一个 block。由于 \(e_1,e_2\) 是任取的，
所有跨边都属于同一个 block。
\end{proof}

\begin{corollary}[对 MPDP 一般图枚举的意义]
任意全局 CCP-Pair 的边界都可以归属到某一个 block 内部。换言之，MPDP 在
一般图上不需要枚举整个 \(S\) 的所有 subset；它只需要在每个 block 内枚举
局部的 \((lb,rb)\)，再用 \textsc{grow} 沿着 block-cut tree 把 block 外
挂接的部分扩展到左侧或右侧。
\end{corollary}

\begin{proof}
对任意全局 CCP-Pair，主命题已经说明其所有跨边属于同一个 block。这个
block 就是全局切分真正发生的位置。block 外部的结构只能通过 cut vertex
挂到这个切口所在 block 上，因此可以由 \textsc{grow} 根据连通性和互斥性
归并到 \(S_L\) 或 \(S_R\)。这正是 MPDP 将全局 subset 枚举替换为
block-local 枚举的正确性基础之一。
\end{proof}

\subsection{grow 如何从 lb/rb 扩展出全局左右侧}

上一小节说明了一个全局 CCP-Pair 的切口一定落在某个 block 内。本小节继续
说明：一旦在这个 block 内选中了局部左右锚点 \(lb\) 和 \(rb\)，为什么
\textsc{grow} 能把它们扩展成当前 DP 状态 \(S\) 上的全局
\((S_{\mathrm{left}},S_{\mathrm{right}})\)。

先把 \textsc{grow} 写成一个普通的图可达性定义。若 \(A\subseteq U\subseteq S\)，
则
\[
\operatorname{grow}(A,U)
=\{v\in U\mid \text{存在 }a\in A\text{，使得 }a\text{ 在 }G[U]\text{ 中可达 }v\}.
\]
也就是说，\(\operatorname{grow}(lb,S\setminus rb)\) 的含义是：从左锚点
\(lb\) 出发，只允许在 \(S\setminus rb\) 里走；所有能走到的点都归到左侧。
这里 \(rb\) 不是普通的未访问集合，而是右侧锚点形成的禁区。

\begin{proposition}[从真实全局 pair 反推 grow]
设 \((S_L^\star,S_R^\star)\) 是当前 connected set \(S\) 上的一个全局
CCP-Pair，且它的所有跨边都落在同一个 block \(B\) 内。令
\[
lb=B\cap S_L^\star,\qquad rb=B\cap S_R^\star.
\]
若 Algorithm 3 正在考察这个 block-local split，则
\[
\operatorname{grow}(lb,S\setminus rb)=S_L^\star,
\qquad
S\setminus \operatorname{grow}(lb,S\setminus rb)=S_R^\star.
\]
\end{proposition}

\begin{proof}
先证 \(S_L^\star\subseteq \operatorname{grow}(lb,S\setminus rb)\)。因为
\((S_L^\star,S_R^\star)\) 是 CCP-Pair，左右之间至少有跨边；这些跨边又都在
\(B\) 内，所以 \(lb\neq\varnothing\)。任取 \(v\in S_L^\star\)。由于
\(G[S_L^\star]\) 连通，存在一条从某个 \(lb\) 中顶点到 \(v\) 的路径，并且
这条路径完全位于 \(S_L^\star\)。而 \(S_L^\star\cap rb=\varnothing\)，所以
这条路径也完全位于 \(S\setminus rb\)。因此 \(v\) 会被 \textsc{grow} 吸入。

再证反向包含。假设存在
\[
v\in \operatorname{grow}(lb,S\setminus rb)\cap S_R^\star.
\]
则存在一条路径 \(P\) 从 \(lb\subseteq S_L^\star\) 走到 \(v\in S_R^\star\)，
并且 \(P\) 完全位于 \(S\setminus rb\)。沿着 \(P\) 从左到右看，必然存在第一条
从 \(S_L^\star\) 跨到 \(S_R^\star\) 的边。这条边是全局 CCP-Pair 的跨边，
根据上一小节主命题，它必须落在 block \(B\) 内。于是它在右侧的端点属于
\(B\cap S_R^\star=rb\)。但 \(P\) 被限制在 \(S\setminus rb\) 中，不能经过
任何 \(rb\) 顶点，矛盾。因此 \(\operatorname{grow}(lb,S\setminus rb)\)
不包含任何 \(S_R^\star\) 中的点，只能等于 \(S_L^\star\)。补集自然就是
\(S_R^\star\)。
\end{proof}

\begin{proposition}[由 block-local split 生成的全局 pair 是 CCP-Pair]
设 \(S\) 是当前 DP 外层循环中的 connected set。Algorithm 3 在某个 block
内枚举出 \(lb\) 和 \(rb\)，并且已经通过检查：
\[
lb\neq\varnothing,\quad rb\neq\varnothing,\quad
G[lb]\text{ 连通},\quad G[rb]\text{ 连通},\quad
lb\text{ 与 }rb\text{ 之间有边}.
\]
令
\[
S_L=\operatorname{grow}(lb,S\setminus rb),\qquad S_R=S\setminus S_L.
\]
则 \((S_L,S_R)\) 是 \(S\) 上的一个 CCP-Pair。
\end{proposition}

\begin{proof}
首先，\(lb\subseteq S_L\)，而 \(rb\cap S_L=\varnothing\)，所以
\(S_L\neq\varnothing\)、\(S_R\neq\varnothing\)，且二者互斥。由于 \(S_L\)
是 \(G[S\setminus rb]\) 中从连通种子 \(lb\) 出发得到的可达闭包，\(G[S_L]\)
连通。

关键是证明 \(G[S_R]\) 也连通。取任意 \(v\in S_R\setminus rb\)。因为
\(G[S]\) 连通，存在一条从 \(v\) 到 \(lb\) 的路径。若某条这样的路径完全不
经过 \(rb\)，则 \(v\) 会在 \(G[S\setminus rb]\) 中从 \(lb\) 可达，从而
\(v\in S_L\)，矛盾。因此从 \(v\) 到 \(lb\) 的任意路径都会遇到 \(rb\)。
取一条简单路径，并令 \(r\) 为从 \(v\) 出发时第一次遇到的 \(rb\) 顶点。
从 \(v\) 到 \(r\) 的前缀不会经过 \(rb\) 的其他顶点，也不能经过 \(S_L\)：
否则 \(v\) 可以先走到某个 \(S_L\) 顶点，再沿 \(S_L\) 内部路径走到 \(lb\)，
这就给出了一条避开 \(rb\) 的可达路径，仍然矛盾。于是该前缀完全位于
\(S_R\)，说明 \(v\) 在 \(G[S_R]\) 中能连到 \(rb\)。

又因为 \(G[rb]\) 连通，所有 \(S_R\) 中的顶点都能连到同一个连通的 \(rb\)
区域，所以 \(G[S_R]\) 连通。最后，Algorithm 3 已经检查 \(lb\) 与 \(rb\)
之间存在边；由于 \(lb\subseteq S_L\)、\(rb\subseteq S_R\)，左右两边之间
存在跨边。因此 \((S_L,S_R)\) 满足 CCP-Pair 的所有条件。
\end{proof}

\begin{corollary}[为什么只需要一次左侧 grow]
Algorithm 3 令
\[
S_L=\operatorname{grow}(lb,S\setminus rb),\qquad S_R=S\setminus S_L,
\]
而不是再单独计算一次 \(\operatorname{grow}(rb,S\setminus lb)\)。原因是：
左侧 \textsc{grow} 已经把所有不需要穿过 \(rb\) 就能从 \(lb\) 到达的点
全部吸走；剩余的点必然都通过 \(rb\) 所在一侧连通起来。因此取补集就得到
合法的 \(S_R\)。
\end{corollary}

\end{document}
